\chapter{GIỚI THIỆU TỔNG QUAN}
\label{ch:gioi_thieu}

\section{Đặt vấn đề}

Xe tự hành và các hệ thống hỗ trợ lái tiên tiến (ADAS) phải hiểu được môi trường giao thông xung
quanh trước khi có thể ra bất kỳ quyết định điều khiển nào. Trong số các thành phần của bài toán
nhận thức đó, phát hiện biển báo giao thông giữ một vị trí đặc biệt: biển báo là kênh giao tiếp
chính thức giữa hạ tầng đường bộ và người lái, và nội dung của nó ràng buộc trực tiếp hành vi hợp
pháp của phương tiện. Một hệ thống bỏ sót biển giới hạn tốc độ hoặc đọc nhầm đèn tín hiệu không
chỉ gây khó chịu mà có thể dẫn tới hậu quả về an toàn và pháp lý.

Về mặt kỹ thuật, đây là bài toán phát hiện đối tượng: với mỗi ảnh đầu vào, mô hình phải đồng thời
\emph{định vị} từng biển báo bằng một hộp bao và \emph{phân loại} nó vào đúng chủng loại. Điều này
khó hơn hẳn bài toán phân lớp ảnh, nơi đầu ra chỉ là một nhãn duy nhất cho cả bức ảnh. Lịch sử
phát triển của lĩnh vực đi qua hai dòng lớn. Dòng hai giai đoạn, tiêu biểu là Faster
R-CNN~\cite{ren2017fasterrcnn}, tách bài toán thành bước đề xuất vùng rồi mới phân loại và tinh
chỉnh hộp bao; độ chính xác cao nhưng chi phí tính toán lớn. Dòng một giai đoạn, khởi đầu từ
YOLO~\cite{redmon2016yolo} và SSD~\cite{liu2016ssd}, dự đoán trực tiếp hộp bao và lớp trên bản đồ
đặc trưng chỉ trong một lượt truyền xuôi, đổi một phần độ chính xác lấy tốc độ đủ cho ứng dụng
thời gian thực.

Năm 2020, DETR~\cite{carion2020detr} mở ra một hướng thứ ba bằng cách áp dụng kiến trúc
transformer~\cite{vaswani2017attention} vào phát hiện đối tượng. DETR loại bỏ hoàn toàn hộp neo
và bước hậu xử lý NMS, thay bằng một tập truy vấn đối tượng học được và phép gán song ánh
Hungarian~\cite{kuhn1955hungarian} giữa dự đoán và nhãn thật. Thiết kế này thanh lịch về mặt khái
niệm, nhưng kèm theo một cái giá ít được nhấn mạnh trong các bài trình bày phổ thông: DETR hội tụ
rất chậm, bản gốc cần khoảng 300 đến 500 chu kỳ huấn luyện trên bộ dữ liệu COCO quy
mô lớn~\cite{lin2014coco}. Đây chính là điểm mà một đồ án học phần với ngân sách GPU hạn chế sẽ
va phải.

Các nút thắt chính mà đề tài phải đối mặt:

\begin{enumerate}[leftmargin=1.5cm]
\item \textbf{Ràng buộc kép giữa độ chính xác và chi phí triển khai.} Mô hình phải chạy được trên
phần cứng nhúng gắn trên xe, nơi dung lượng bộ nhớ và ngân sách năng lượng đều bị giới hạn, nhưng
vẫn phải đủ chính xác để không bỏ sót biển báo.

\item \textbf{Dữ liệu quy mô vừa và mất cân bằng.} Bộ dữ liệu sử dụng chỉ có 6.012 hộp bao cho 15
lớp, với tỉ số mất cân bằng giữa lớp phổ biến nhất và lớp hiếm nhất lên tới 35,77 lần. Mất cân
bằng lớp là một vấn đề đã được ghi nhận rộng rãi trong học sâu~\cite{johnson2019imbalance} và ảnh
hưởng trực tiếp tới các lớp thiểu số.

\item \textbf{Vật thể nhỏ.} Biển báo ở xa chỉ chiếm vài chục điểm ảnh trong khung hình. Phát hiện
vật thể nhỏ là điểm yếu cố hữu của các bộ phát hiện hiện đại~\cite{tong2020smallobject}, và mức
độ nghiêm trọng của nó phụ thuộc mạnh vào phân bố kích thước trong dữ liệu huấn luyện.

\item \textbf{So sánh công bằng giữa hai hệ sinh thái khác nhau.} YOLO và DETR đến từ hai thư viện
khác nhau, dùng hai định dạng nhãn khác nhau và hai cách tính mAP khác nhau. Nếu không kiểm soát,
phần chênh lệch do công cụ đo sẽ lẫn vào phần chênh lệch do mô hình.
\end{enumerate}

\noindent\textbf{Câu hỏi nghiên cứu 1.1} \textit{Với một bộ dữ liệu biển báo quy mô vừa và ngân
sách huấn luyện hạn chế, kiến trúc một giai đoạn dựa trên tích chập và kiến trúc transformer dự
đoán theo tập hợp cho kết quả khác nhau như thế nào về độ chính xác, tốc độ và dung lượng?}

\noindent\textbf{Câu hỏi nghiên cứu 1.2} \textit{Có kiến trúc nhẹ nào thay thế được mô hình nền
mà không đánh đổi độ chính xác, và các kỹ thuật nén mô hình như lượng tử hoá hay chưng cất tri
thức mang lại lợi ích gì trong bối cảnh này?}

\noindent\textbf{Câu hỏi nghiên cứu 1.3} \textit{Chỉ số mAP đo trên tập kiểm tra phản ánh đến mức
nào năng lực thực tế của mô hình khi gặp cảnh đường phố thật, và đặc điểm nào của bộ dữ liệu giải
thích khoảng cách đó?}

\section{Mục tiêu đề tài}

Đề tài hướng đến các mục tiêu chính sau:

\begin{itemize}[leftmargin=1.5cm]
\item Xây dựng một quy trình xử lý hoàn chỉnh và tái lập được, đi từ dữ liệu thô tới mô hình đã
huấn luyện và bảng số liệu đánh giá, trong đó mọi bước đều tồn tại dưới dạng mô-đun độc lập chạy
được từ dòng lệnh và có kiểm thử tự động.

\item Thiết lập một mô hình nền đáng tin cậy bằng YOLOv8n và một mô hình đối chứng bằng
DETR-ResNet50, đo trên cùng một tập kiểm tra với cùng một hạt giống ngẫu nhiên.

\item Khảo sát ba kiến trúc nhẹ thay thế và hai nhóm kỹ thuật nén mô hình, nhằm trả lời câu hỏi
liệu có thể giảm chi phí triển khai mà không mất độ chính xác.

\item Phân tích khám phá dữ liệu ở mức đủ sâu để giải thích được kết quả, thay vì chỉ báo cáo con
số cuối cùng.

\item Báo cáo trung thực cả những kết quả đi ngược giả thuyết ban đầu, kèm theo phân tích nguyên
nhân.
\end{itemize}

\section{Đối tượng và phạm vi nghiên cứu}

\textbf{Đối tượng nghiên cứu:} các mô hình học sâu phát hiện đối tượng áp dụng cho biển báo giao
thông, xét đồng thời trên ba trục độ chính xác, tốc độ suy luận và dung lượng mô hình.

\textbf{Phạm vi nghiên cứu:} đề tài giới hạn ở bộ dữ liệu \textit{pkdarabi/cardetection} với 15
lớp gốc, giữ nguyên không ánh xạ lại nhãn. Việc huấn luyện được thực hiện trên môi trường Google
Colab với GPU đơn, ngân sách 30 chu kỳ cho nhánh YOLO và 10 chu kỳ cho nhánh DETR. Đề tài
\emph{không} bao gồm: thu thập và gán nhãn dữ liệu mới, tinh chỉnh siêu tham số bằng tìm kiếm tự
động, huấn luyện DETR tới hội tụ hoàn toàn, và triển khai thực tế trên phần cứng nhúng. Những
giới hạn này được nêu rõ ngay từ đầu vì chúng ảnh hưởng trực tiếp tới cách diễn giải kết quả ở
Chương~\ref{ch:ket_qua}.

\section{Đóng góp của đề tài}

\begin{enumerate}[leftmargin=1.5cm]
\item Một quy trình tái lập được cho bài toán phát hiện biển báo, trong đó toàn bộ bộ sinh số ngẫu
nhiên được gieo hạt tập trung, mọi đường dẫn được quản lý ở một nơi duy nhất, và bộ kiểm thử tự
động gồm 43 trường hợp chạy trên CPU trong khoảng một giây.

\item Một bảng so sánh bốn nhóm thực nghiệm trên cùng một tập kiểm tra: mô hình nền, kiến trúc
nhẹ, lượng tử hoá và chưng cất tri thức --- cho phép đọc được sự đánh đổi giữa ba trục thay vì chỉ
một con số mAP.

\item Một phân tích thành phần bộ dữ liệu chỉ ra rằng 38,7\% số ảnh là ảnh cắt cận cảnh kiểu
GTSRB~\cite{stallkamp2012gtsrb}, khiến phân bố kích thước vật thể trở nên lưỡng cực. Phát hiện này
giải thích được khoảng cách giữa chỉ số trên tập kiểm tra và hành vi quan sát được trên video
đường phố thực tế, và nó liên hệ trực tiếp tới hiện tượng thiên lệch bộ dữ liệu đã được mô tả
trong tài liệu~\cite{torralba2011bias}.

\item Hai kết quả âm được báo cáo đầy đủ thay vì lược bỏ: chưng cất tri thức làm giảm độ chính xác
so với đối chứng, và lượng tử hoá INT8 làm giảm tốc độ thay vì tăng trên môi trường thử nghiệm.
\end{enumerate}

\section{Bố cục báo cáo}

Chương~\ref{ch:gioi_thieu} giới thiệu tổng quan về bài toán, mục tiêu và phạm vi.
Chương~\ref{ch:co_so} trình bày cơ sở lý thuyết về các họ kiến trúc phát hiện đối tượng, các hàm
mất mát và hệ thống độ đo. Chương~\ref{ch:phuong_phap} mô tả bộ dữ liệu, kết quả phân tích khám
phá, kiến trúc quy trình và thiết kế của bốn nhóm thực nghiệm. Chương~\ref{ch:ket_qua} trình bày
và bàn luận kết quả, bao gồm phần phân tích lỗi và các kết quả trái với kỳ vọng.
Chương~\ref{ch:tong_ket} tổng kết, nêu hạn chế và hướng phát triển.
